\begin{abstract}
Accurate parameterization of reaction rate laws is essential for reactor design,  process optimization, and scale-up in chemical engineering applications. The  acid-catalyzed iodination of acetone was studied to determine its rate law and temperature dependence using the method of initial rates. Reaction orders were determined using linear, multiple linear, and nonlinear regression. Nonlinear regression yielded the most precise parameter estimates, with 95\% confidence intervals of $\alpha = 1.026 \pm 0.063$, $\beta = 0.954 \pm 0.068$, and $\gamma = -0.134 \pm 0.67$ for acetone, hydrochloric acid (HCl), and iodine, respectively, consistent with a pseudo-equilibrium mechanism. The activation energy and pre-exponential factor were determined from a linearized Arrhenius plot to be $E_a = (114.4 \pm 5.6)$ kJ/mol and $\ln(A) = (36.0 \pm 2.2)$ (standard error). The measured activation energy is in statistical agreement with the literature consensus of approximately 86--89 kJ/mol. However, this agreement is tenuous given that only three temperature points were used, producing a critical t-value of ~12.7 and a confidence interval wide enough to encompass nearly any physically reasonable activation energy.  The experimentally determined rate law was applied to numerically model autocatalytic behavior in a BSTR and to size a CSTR for iodoacetone production at 353 K and 80\% acetone conversion, yielding a nominal volume of 274 L. A Monte Carlo simulation ($N = 100{,}000$) produced a 90\% confidence interval of 58--1,290 L, with uncertainty dominated by the Arrhenius parameters $E_a$ and $\ln(A)$. Additional rate measurements across a wider temperature range between 303--353 K are recommended to most efficiently constrain reactor volume uncertainty along with having more statistical power compare the apparent activation energy to literature values.
\vspace{1em}

A practical survey of pneumatic pressure measurement instruments was conducted to understand their functional roles in process systems. Instruments examined included Bourdon gauges, pressure transducers, mass flow controllers, and control valves. Operation of a pneumatic pressure rack enabled direct observation of pressure regulation, signal transmission, and feedback control behavior. Differences between mechanical and electronic measurement devices were evaluated in terms of response time, sensitivity, and stability. Control valve performance under varying setpoints illustrated the importance of proper controller tuning in maintaining steady-state operation. This exercise provided practical insight into the integration of instrumentation and control hardware and emphasized their role in ensuring safe, stable, and efficient industrial processes.


\end{abstract}